% !TEX program = pdflatex
% ==========================================================
% Universidad Santo Tomás
% Fundamentos de Matemática Básica
% Clase 3: Potencias
% Versión Premium
% Luis Tulio González Alcaino
% ==========================================================

\newif\ifhandout
%\handoutfalse
\handouttrue

\ifhandout
\documentclass[aspectratio=169,10pt,handout]{beamer}
\else
\documentclass[aspectratio=169,10pt]{beamer}
\fi

% ----------------------------------------------------------
% PAQUETES
% ----------------------------------------------------------
\usepackage[spanish,es-noshorthands]{babel}
\usepackage[T1]{fontenc}
\usepackage[utf8]{inputenc}

\usepackage{lmodern}
\usepackage{amsmath,amssymb,amsfonts}
\usepackage{ragged2e}
\usepackage{enumitem}
\usepackage{xcolor}
\usepackage{tikz}
\usetikzlibrary{positioning,calc}
\usepackage{fontawesome5}

% ----------------------------------------------------------
% COLORES INSTITUCIONALES
% ----------------------------------------------------------
\definecolor{USTBlue}{RGB}{0,70,127}
\definecolor{USTBlueDark}{RGB}{0,45,85}
\definecolor{USTTeal}{RGB}{0,153,117}
\definecolor{USTGray}{RGB}{120,120,120}
\definecolor{USTSoft}{RGB}{248,249,250}
\definecolor{USTText}{RGB}{35,40,45}
\definecolor{USTLightBlue}{RGB}{227,239,248}
\definecolor{USTLightTeal}{RGB}{226,245,240}
\definecolor{USTRedSoft}{RGB}{252,235,235}

% ----------------------------------------------------------
% TEMA
% ----------------------------------------------------------
\usetheme{Madrid}
\usecolortheme[named=USTBlue]{structure}
\setbeamertemplate{navigation symbols}{}

% ----------------------------------------------------------
% AJUSTES VISUALES
% ----------------------------------------------------------
\setlength{\parskip}{0.2em}
\setlist[itemize]{itemsep=0.35em}
\setlist[enumerate]{itemsep=0.35em}

\setbeamercolor{normal text}{fg=USTText,bg=white}
\setbeamercolor{frametitle}{fg=white,bg=USTBlueDark}
\setbeamercolor{title}{fg=USTBlueDark}
\setbeamercolor{subtitle}{fg=USTTeal}

\setbeamercolor{block title}{bg=USTLightBlue,fg=USTBlueDark}
\setbeamercolor{block body}{bg=USTSoft,fg=black}

\setbeamercolor{block title example}{bg=USTLightTeal,fg=USTBlueDark}
\setbeamercolor{block body example}{bg=USTSoft,fg=black}

\setbeamercolor{block title alerted}{bg=USTRedSoft,fg=black}
\setbeamercolor{block body alerted}{bg=red!2,fg=black}

% ----------------------------------------------------------
% FRAMETITLE
% ----------------------------------------------------------
\setbeamertemplate{frametitle}{
	\vspace*{-0.03cm}
	\begin{beamercolorbox}[wd=\paperwidth,ht=0.95cm,dp=0.2cm,leftskip=0.45cm,rightskip=0.25cm]{frametitle}
		{\large\bfseries \insertframetitle}
	\end{beamercolorbox}
}

% ----------------------------------------------------------
% PIE DE PÁGINA
% ----------------------------------------------------------
\setbeamercolor{author in head/foot}{bg=USTSoft,fg=USTGray}
\setbeamercolor{date in head/foot}{bg=USTSoft,fg=USTGray}
\setbeamercolor{title in head/foot}{bg=USTSoft,fg=USTGray}

\setbeamertemplate{footline}{%
	\leavevmode%
	\hbox{%
		\begin{beamercolorbox}[wd=.82\paperwidth,ht=2.8ex,dp=1ex,leftskip=0.4cm]{author in head/foot}%
			\scriptsize
			\textbf{Fundamentos de Matemática Básica}
			\quad|\quad
			Agronomía
			\quad|\quad
			Universidad Santo Tomás Talca
		\end{beamercolorbox}%
		\begin{beamercolorbox}[wd=.18\paperwidth,ht=2.8ex,dp=1ex,rightskip=0.35cm plus1fil]{date in head/foot}%
			\scriptsize \insertframenumber/\inserttotalframenumber
		\end{beamercolorbox}%
	}%
}

% ----------------------------------------------------------
% COMANDOS
% ----------------------------------------------------------
\newcommand{\destacar}[1]{\textcolor{USTBlueDark}{\textbf{#1}}}
\newcommand{\alerta}[1]{\textcolor{red!70!black}{\textbf{#1}}}
\newcommand{\pp}{\pause}

\newcommand{\iconoInicio}{\faPlayCircle}
\newcommand{\iconoDesarrollo}{\faBrain}
\newcommand{\iconoPractica}{\faPenSquare}
\newcommand{\iconoAgro}{\faSeedling}
\newcommand{\iconoCierre}{\faCheckCircle}
\newcommand{\iconoIdea}{\faLightbulb}
\newcommand{\iconoError}{\faExclamationTriangle}

\newenvironment{metaBox}
{\begin{block}{\faBullseye\ \ Meta de aprendizaje}}
	{\end{block}}

\newenvironment{ideaBox}
{\begin{exampleblock}{\iconoIdea\ \ Idea clave}}
	{\end{exampleblock}}

% ----------------------------------------------------------
% PORTADA PREMIUM
% ----------------------------------------------------------
\newcommand{\PortadaProfesionalUST}{%
	\begin{frame}[plain]
		\begin{tikzpicture}[remember picture,overlay]
			
			\fill[white] (current page.south west) rectangle (current page.north east);
			\fill[USTBlueDark] (current page.north west) rectangle ([yshift=-1.4cm]current page.north east);
			\fill[USTTeal] ([yshift=-1.4cm]current page.north west) rectangle ([yshift=-1.58cm]current page.north east);
			\fill[USTSoft] (current page.south west) rectangle ([xshift=2.2cm]current page.north west);
			
			\fill[USTBlue!8] ([xshift=1.1cm,yshift=-2.0cm]current page.north west) circle (0.72cm);
			\fill[USTTeal!10] ([xshift=1.7cm,yshift=-4.2cm]current page.north west) circle (0.45cm);
			\fill[USTBlue!6] ([xshift=-1.1cm,yshift=1.0cm]current page.south east) circle (1.05cm);
			
			\fill[black!6,rounded corners=8pt]
			([xshift=-5.0cm,yshift=-3.1cm]current page.center) rectangle
			([xshift=5.0cm,yshift=2.7cm]current page.center);
			
			\fill[white,rounded corners=8pt]
			([xshift=-5.15cm,yshift=-2.95cm]current page.center) rectangle
			([xshift=4.85cm,yshift=2.85cm]current page.center);
			
			\fill[USTTeal] ([xshift=-5.15cm,yshift=-2.95cm]current page.center) rectangle
			([xshift=-4.85cm,yshift=2.85cm]current page.center);
			
			\node[anchor=north east,text=white]
			at ([xshift=-0.7cm,yshift=-0.5cm]current page.north east)
			{\small\bfseries Universidad Santo Tomás \quad Ciencias Básicas Talca};
			
			\node[align=center] at ([xshift=0.15cm,yshift=-0.05cm]current page.center){%
				\begin{minipage}{0.68\paperwidth}
					\centering
					{\Large\bfseries\color{USTText}\insertauthor\par}
					\vspace{0.7cm}
					{\fontsize{24}{28}\selectfont\bfseries\color{USTBlueDark}\inserttitle\par}
					\vspace{0.25cm}
					{\large\bfseries\color{USTTeal}\insertsubtitle\par}
					\vspace{0.55cm}
					{\normalsize\color{USTGray}\faSeedling\ Agronomía \quad \faUniversity\ Universidad Santo Tomás\par}
					\vspace{0.65cm}
					{\small\color{USTGray}\insertdate\par}
				\end{minipage}
			};
			
			\draw[USTGray!35,line width=0.35pt]
			([xshift=0.9cm,yshift=0.72cm]current page.south west) --
			([xshift=-0.9cm,yshift=0.72cm]current page.south east);
			
			\node[anchor=south,text=USTGray]
			at ([yshift=0.2cm]current page.south)
			{\scriptsize
				\textbf{Fundamentos de Matemática Básica}
				\;|\;
				Clase Premium
				\;|\;
				Potencias};
		\end{tikzpicture}
	\end{frame}
}

% ----------------------------------------------------------
% PORTADILLA DE SECCIÓN
% ----------------------------------------------------------
\newcommand{\SectionFrame}[3]{%
	\begin{frame}[plain]
		\begin{tikzpicture}[remember picture,overlay]
			\fill[USTBlueDark] (current page.south west) rectangle (current page.north east);
			\fill[USTTeal] (current page.south west) rectangle ([yshift=0.23cm]current page.north east);
			\fill[white,opacity=0.08] ([xshift=-1cm,yshift=1cm]current page.south east) circle (1.8cm);
			\fill[white,opacity=0.06] ([xshift=1.3cm,yshift=-0.8cm]current page.north west) circle (1.25cm);
			
			\node[anchor=center,text=white] at ([yshift=0.9cm]current page.center)
			{{\fontsize{34}{36}\selectfont #1}};
			
			\node[anchor=center,text=white] at (current page.center)
			{{\LARGE\bfseries #2}};
			
			\node[anchor=center,text=white!90] at ([yshift=-1.0cm]current page.center)
			{{\large #3}};
		\end{tikzpicture}
	\end{frame}
}

% ----------------------------------------------------------
% DATOS DEL CURSO
% ----------------------------------------------------------
\title{Fundamentos de Matemática Básica}
\subtitle{Clase 3: Potencias}
\author{Prof. Luis González Alcaino}
\date{Universidad Santo Tomás -- Agronomía 2026}

% ==========================================================
% DOCUMENTO
% ==========================================================
\begin{document}
	
	\PortadaProfesionalUST
	
	% ----------------------------------------------------------
	\begin{frame}{Panorama de la clase}
		\begin{metaBox}
			Comprender el concepto de potencia, aplicar sus propiedades y resolver situaciones contextualizadas en Agronomía.
		\end{metaBox}
		
		\pp
		
		\begin{block}{\faListOl\ \ Ruta de trabajo}
			\begin{enumerate}
				\item Concepto de potencia
				\item Aplicaciones en Agronomía
				\item Propiedades fundamentales
				\item Patrones numéricos
				\item Ejercicios y problemas contextualizados
				\item Síntesis final
			\end{enumerate}
		\end{block}
	\end{frame}
	
	% ----------------------------------------------------------
	\SectionFrame{\iconoInicio}{Inicio}{Comprender el concepto de potencia}
	
	% ----------------------------------------------------------
	\begin{frame}{\iconoInicio\ \ Concepto de potencia}
		\begin{block}{Definición}
			Una \destacar{potencia} representa una multiplicación repetida de un mismo número:
			\[
			a^n = a\cdot a\cdot a \cdots a
			\]
			donde:
			\begin{itemize}
				\item \(a\) es la \destacar{base}
				\item \(n\) es el \destacar{exponente}
			\end{itemize}
		\end{block}
		
		\pause
		
		\begin{exampleblock}{Ejemplos}
			\[
			3^4 = 3\cdot3\cdot3\cdot3 = 81
			\qquad
			2^5 = 32
			\]
			\[
			10^3 = 1000
			\]
		\end{exampleblock}
	\end{frame}
	
	% ----------------------------------------------------------
	\begin{frame}{\iconoInicio\ \ Otra forma de interpretar una potencia}
		\begin{block}{Definición ampliada}
			Una \destacar{potencia} puede escribirse como:
			\[
			a^n=\underbrace{a\cdot a\cdot a\cdots a}_{n \text{ veces}}
			\]
		\end{block}
		
		\pp
		
		\begin{exampleblock}{Ejemplos}
			\[
			2^4 = 2\cdot2\cdot2\cdot2 = 16
			\qquad
			5^2 = 25
			\]
			\[
			10^{-2}=\frac{1}{10^2}=\frac{1}{100}=0{,}01
			\]
		\end{exampleblock}
	\end{frame}
	
	% ----------------------------------------------------------
	\SectionFrame{\iconoAgro}{Contexto disciplinar}{¿Por qué estudiar potencias en Agronomía?}
	
	% ----------------------------------------------------------
	\begin{frame}{\iconoAgro\ \ Aplicaciones en Agronomía}
		\begin{columns}[T,totalwidth=\textwidth]
			\begin{column}{0.52\textwidth}
				\begin{block}{Situaciones reales}
					\begin{itemize}
						\item Número de semillas por hectárea
						\item Población de bacterias en suelo
						\item Concentración de micronutrientes
						\item Conversión de unidades y escalas
					\end{itemize}
				\end{block}
			\end{column}
			
			\begin{column}{0.44\textwidth}
				\begin{exampleblock}{Ejemplo rápido}
					\[
					0{,}00045\ \text{kg}=4{,}5\times10^{-4}\ \text{kg}
					\]
					La segunda forma es más clara para calcular, comparar y comunicar.
				\end{exampleblock}
			\end{column}
		\end{columns}
		
		\vspace{0.2cm}
		\centering
		{\small \textcolor{USTGray}{La matemática permite simplificar, modelar y validar información técnica.}}
	\end{frame}
	
	% ----------------------------------------------------------
	\SectionFrame{\iconoDesarrollo}{Desarrollo}{Propiedades de las potencias}
	
	% ----------------------------------------------------------
	\begin{frame}{\iconoDesarrollo\ \ Propiedades I}
		\begin{block}{1. Producto de potencias de igual base}
			\[
			a^m \cdot a^n = a^{m+n}
			\]
		\end{block}
		
		\begin{exampleblock}{Ejemplo}
			\[
			2^3\cdot2^4 = 2^{3+4}=2^7=128
			\]
		\end{exampleblock}
		
		\pause
		
		\begin{block}{2. Cociente de potencias de igual base}
			\[
			\frac{a^m}{a^n}=a^{m-n}, \qquad a\neq 0
			\]
		\end{block}
		
		\begin{exampleblock}{Ejemplo}
			\[
			\frac{5^6}{5^2}=5^{6-2}=5^4=625
			\]
		\end{exampleblock}
	\end{frame}
	
	% ----------------------------------------------------------
	\begin{frame}{\iconoDesarrollo\ \ Propiedades II}
		\begin{block}{3. Potencia de una potencia}
			\[
			(a^m)^n = a^{m\cdot n}
			\]
		\end{block}
		
		\begin{exampleblock}{Ejemplo}
			\[
			(3^2)^4 = 3^{2\cdot4}=3^8=6561
			\]
		\end{exampleblock}
		
		\pause
		
		\begin{block}{4. Potencia de un producto}
			\[
			(ab)^n = a^n b^n
			\]
		\end{block}
		
		\begin{exampleblock}{Ejemplo}
			\[
			(2\cdot5)^3 = 2^3\cdot5^3 = 8\cdot125 = 1000
			\]
		\end{exampleblock}
	\end{frame}
	
	% ----------------------------------------------------------
	\begin{frame}{\iconoDesarrollo\ \ Propiedades III}
		\begin{block}{5. Potencia de un cociente}
			\[
			\left(\frac{a}{b}\right)^n = \frac{a^n}{b^n}, \qquad b\neq 0
			\]
		\end{block}
		
		\begin{exampleblock}{Ejemplo}
			\[
			\left(\frac{2}{3}\right)^3 = \frac{2^3}{3^3}=\frac{8}{27}
			\]
		\end{exampleblock}
		
		\pause
		
		\begin{block}{6. Exponente cero}
			\[
			a^0 = 1, \qquad a\neq 0
			\]
		\end{block}
		
		\begin{exampleblock}{Ejemplo}
			\[
			7^0 = 1
			\]
		\end{exampleblock}
	\end{frame}
	
	% ----------------------------------------------------------
	\begin{frame}{\iconoDesarrollo\ \ Propiedades IV}
		\begin{block}{7. Exponente uno}
			\[
			a^1 = a
			\]
		\end{block}
		
		\begin{exampleblock}{Ejemplo}
			\[
			9^1 = 9
			\]
		\end{exampleblock}
		
		\pause
		
		\begin{block}{8. Exponente negativo}
			\[
			a^{-n}=\frac{1}{a^n}, \qquad a\neq 0
			\]
		\end{block}
		
		\begin{exampleblock}{Ejemplo}
			\[
			2^{-3}=\frac{1}{2^3}=\frac{1}{8}
			\]
		\end{exampleblock}
	\end{frame}
	
	% ----------------------------------------------------------
	\begin{frame}{\iconoDesarrollo\ \ Propiedades V}
		\begin{block}{9. Base igual a uno}
			\[
			1^n = 1
			\]
		\end{block}
		
		\begin{exampleblock}{Ejemplo}
			\[
			1^8 = 1
			\]
		\end{exampleblock}
		
		\pause
		
		\begin{block}{10. Base igual a cero}
			\[
			0^n = 0, \qquad n>0
			\]
		\end{block}
		
		\begin{exampleblock}{Ejemplo}
			\[
			0^5 = 0
			\]
		\end{exampleblock}
	\end{frame}
	
	% ----------------------------------------------------------
	\begin{frame}{\iconoError\ \ Error frecuente}
		\begin{alertblock}{Cuidado}
			No se pueden sumar ni restar exponentes en una suma de potencias:
			\[
			a^m + a^n \neq a^{m+n}
			\]
		\end{alertblock}
		
		\pause
		
		\begin{exampleblock}{Ejemplo}
			\[
			2^2 + 2^3 = 4 + 8 = 12
			\]
			pero
			\[
			2^{2+3}=2^5=32
			\]
			Por lo tanto,
			\[
			2^2 + 2^3 \neq 2^5
			\]
		\end{exampleblock}
	\end{frame}
	
	% ----------------------------------------------------------
	\SectionFrame{\iconoDesarrollo}{Patrones}{Regularidades numéricas}
	
	% ----------------------------------------------------------
	\begin{frame}{\iconoDesarrollo\ \ Patrones}
		Observe la secuencia:
		\[
		2^1 = 2
		\qquad
		2^2 = 4
		\qquad
		2^3 = 8
		\qquad
		2^4 = 16
		\]
		
		\pause
		
		\begin{block}{Preguntas}
			\begin{itemize}
				\item ¿Qué ocurre cuando el exponente aumenta en 1?
				\item ¿Cómo crecerá la secuencia?
			\end{itemize}
		\end{block}
	\end{frame}
	
	% ----------------------------------------------------------
	\begin{frame}{\iconoDesarrollo\ \ Potencias de 10}
		\begin{columns}[T,totalwidth=\textwidth]
			\begin{column}{0.48\textwidth}
				\begin{block}{Exponente positivo}
					\[
					10^1 = 10
					\]
					\[
					10^2 = 100
					\]
					\[
					10^3 = 1000
					\]
					\[
					10^4 = 10000
					\]
				\end{block}
			\end{column}
			
			\begin{column}{0.48\textwidth}
				\begin{exampleblock}{Exponente negativo}
					\[
					10^0 = 1
					\]
					\[
					10^{-1} = 0{,}1
					\]
					\[
					10^{-2} = 0{,}01
					\]
					\[
					10^{-3} = 0{,}001
					\]
				\end{exampleblock}
			\end{column}
		\end{columns}
		
		\vspace{0.2cm}
		\begin{ideaBox}
			Los exponentes permiten describir regularidades de crecimiento o disminución en escala decimal.
		\end{ideaBox}
	\end{frame}
	
	% ----------------------------------------------------------
	\SectionFrame{\iconoPractica}{Práctica}{Ejercicios y problemas contextualizados}
	
	% ----------------------------------------------------------
	\begin{frame}{\iconoPractica\ \ Ejercicios}
		Resuelva:
		\begin{enumerate}
			\item \(2^4 \cdot 2^3\)
			\item \((3^2)^3\)
			\item \(5^3 \cdot 5^2\)
			\item \(\dfrac{10^6}{10^2}\)
			\item En un cultivo se plantan \(3^4\) semillas por parcela. Si hay \(3^2\) parcelas, ¿cuántas semillas hay en total?
		\end{enumerate}
	\end{frame}
	
	% ----------------------------------------------------------
	\begin{frame}{\iconoAgro\ \ Ejercicio contextualizado}
		\begin{block}{Situación agrícola}
			Una parcela experimental utiliza:
			\[
			2^5
			\]
			semillas por metro cuadrado.
			
			Si el cultivo cubre:
			\[
			2^3
			\]
			metros cuadrados, ¿cuántas semillas se utilizan en total?
		\end{block}
		
		\pause
		
		\begin{exampleblock}{Resolución}
			\[
			2^5 \cdot 2^3 = 2^{5+3} = 2^8
			\]
			\[
			2^8 = 256
			\]
			Se utilizan aproximadamente \textbf{256 semillas}.
		\end{exampleblock}
	\end{frame}
	
	% ----------------------------------------------------------
	\begin{frame}{\iconoAgro\ \ Crecimiento de bacterias en el suelo}
		\begin{block}{Contexto}
			Una colonia bacteriana se duplica cada hora.
			
			La cantidad inicial es:
			\[
			2^3
			\]
			bacterias.
			
			Después de 4 horas:
			\[
			2^3 \cdot 2^4
			\]
		\end{block}
		
		\pause
		
		\begin{exampleblock}{Cálculo}
			\[
			2^3 \cdot 2^4 = 2^{7}
			\]
			\[
			2^7 = 128
			\]
			Después de 4 horas hay \textbf{128 bacterias}.
		\end{exampleblock}
	\end{frame}
	
	% ----------------------------------------------------------
	\begin{frame}{\iconoAgro\ \ Fertilización}
		\begin{block}{Situación}
			Un fertilizante se aplica en dosis proporcionales a potencias de 10.
			
			En un ensayo se utiliza:
			\[
			5\times10^2
			\]
			gramos por hectárea.
			
			Si el cultivo ocupa:
			\[
			10^3
			\]
			hectáreas experimentales, calcule la cantidad total.
		\end{block}
		
		\pause
		
		\begin{exampleblock}{Resolución}
			\[
			(5\times10^2)(10^3)=5\times10^{5}
			\]
			\[
			500000\text{ gramos}
			\]
		\end{exampleblock}
	\end{frame}
	
	% ----------------------------------------------------------
	\begin{frame}{\iconoPractica\ \ Ejercicios contextualizados}
		\begin{enumerate}
			\item En una parcela experimental se plantan \(2^6\) semillas de trigo por sector. Si el terreno tiene \(2^3\) sectores iguales, ¿cuántas semillas se plantan en total?
			
			\pause
			
			\item Una bacteria del suelo se duplica cada hora. Si inicialmente hay \(2^2\) bacterias, ¿cuántas habrá después de 6 horas?
			
			\pause
			
			\item Un cultivo utiliza una densidad de \(5^3\) plantas por unidad experimental. Si se instalan \(5^2\) unidades experimentales, ¿cuántas plantas se utilizan en total?
			
			\pause
			
			\item En un laboratorio agrícola se almacenan muestras de suelo. Cada caja contiene \(3^4\) muestras. Si se tienen \(3^2\) cajas, ¿cuántas muestras hay en total?
		\end{enumerate}
	\end{frame}
	
	% ----------------------------------------------------------
	\SectionFrame{\iconoCierre}{Cierre}{Síntesis de la sesión}
	
	% ----------------------------------------------------------
	\begin{frame}{\iconoCierre\ \ Síntesis}
		\begin{block}{Ideas clave}
			\begin{itemize}
				\item Las potencias representan multiplicaciones repetidas.
				\item Las propiedades de las potencias permiten simplificar cálculos.
				\item Las potencias aparecen frecuentemente en fenómenos de crecimiento.
				\item En Agronomía permiten modelar poblaciones, concentraciones y escalas.
			\end{itemize}
		\end{block}
		
		\pause
		
		\begin{ideaBox}
			Aprender potencias ayuda a calcular mejor, interpretar fenómenos reales y preparar el paso hacia la notación científica.
		\end{ideaBox}
	\end{frame}
	
\end{document}